\documentclass[tikz,border=10pt]{standalone}
\usepackage[no-math]{fontspec}
\usepackage[brazil]{babel}
\usepackage{amsmath,amssymb}
\IfFileExists{C:/Windows/Fonts/Lato_Regular.ttf}{%
  \setmainfont{Lato}[
    Path = C:/Windows/Fonts/,
    UprightFont = Lato_Regular.ttf,
    BoldFont = Lato_Bold.ttf,
    ItalicFont = Lato_Italic.ttf,
    BoldItalicFont = Lato_BoldItalic.ttf
  ]
  \setsansfont{Lato}[
    Path = C:/Windows/Fonts/,
    UprightFont = Lato_Regular.ttf,
    BoldFont = Lato_Bold.ttf,
    ItalicFont = Lato_Italic.ttf,
    BoldItalicFont = Lato_BoldItalic.ttf
  ]
}{%
  \RequirePackage[default]{lato}
}
\renewcommand{\familydefault}{\sfdefault}

\usepackage{xcolor}
\usepackage{tikz}
\usepackage{pgfplots}
\pgfplotsset{compat=1.18}
\usetikzlibrary{arrows.meta,positioning,shapes.geometric,fit,backgrounds,calc,decorations.pathreplacing,patterns}

% Paleta institucional UEFS / Bioengenharia de Solos
\definecolor{bsNavy}{HTML}{1F4268}
\definecolor{bsNavyDark}{HTML}{152E4A}
\definecolor{bsGold}{HTML}{B38A3D}
\definecolor{bsGoldDark}{HTML}{8C6927}
\definecolor{bsTeal}{HTML}{0E7C7B}
\definecolor{bsForest}{HTML}{2D6A4F}
\definecolor{bsAlert}{HTML}{9E2A2B}
\definecolor{bsEarth}{HTML}{8C5A3C}
\definecolor{bsWater}{HTML}{2E86AB}
\definecolor{bsInk}{HTML}{26313D}
\definecolor{bsMuted}{HTML}{5B6573}

\definecolor{bgNavy}{HTML}{EDF2F7}
\definecolor{bgGold}{HTML}{F8F3E9}
\definecolor{bgTeal}{HTML}{EBF6F5}
\definecolor{bgForest}{HTML}{EEF5F0}
\definecolor{bgAlert}{HTML}{FDF2F2}

\tikzset{
  >=Stealth,
  bsBox/.style={
    draw=#1, fill=white, rounded corners=5pt, line width=1.15pt,
    inner sep=7pt, align=center, font=\footnotesize\sffamily, text=bsInk
  },
  bsTint/.style={
    draw=#1, fill=#1!8!white, rounded corners=5pt, line width=1.15pt,
    inner sep=7pt, align=center, font=\footnotesize\sffamily, text=bsInk
  },
  bsBadge/.style={
    fill=#1, text=white, rounded corners=9pt, inner xsep=8pt, inner ysep=3pt,
    font=\scriptsize\bfseries\sffamily
  },
  bsTitle/.style={font=\normalsize\bfseries\sffamily, text=bsNavy, align=center},
  bsSub/.style={font=\footnotesize\sffamily, text=bsMuted, align=center},
  bsArrow/.style={->, >=Stealth, line width=1.35pt},
  bsDArrow/.style={<->, >=Stealth, line width=1.35pt}
}

\tikzset{
  pill/.style={bsBadge=#1, minimum height=0pt, minimum width=0pt, text width=}
}

\begin{document}
\begin{tikzpicture}[node distance=0.65cm]

% 5 cartões sequenciais equilibrados
\node[bsTint=bsNavy, text width=3.05cm, minimum height=2.75cm] (p1) at (0, 0) {%
  \tikz\node[pill=bsNavy]{ETAPA 1};\\[6pt]
  \textbf{Regularização \&\\Remoção de Torrões}\\[4pt]
  {\scriptsize Nivelamento da superfície\\para eliminar vazios\\e evitar efeito tenda}
};

\node[bsTint=bsTeal, text width=3.05cm, minimum height=2.75cm, right=of p1] (p2) {%
  \tikz\node[pill=bsTeal]{ETAPA 2};\\[6pt]
  \textbf{Semeadura \&\\Adubação de Base}\\[4pt]
  {\scriptsize Distribuição uniforme de\\sementes e corretivos\\antes da cobertura}
};

\node[bsTint=bsEarth, text width=3.05cm, minimum height=2.75cm, right=of p2] (p3) {%
  \tikz\node[pill=bsEarth]{ETAPA 3};\\[6pt]
  \textbf{Valeta de Ancoragem\\na Crista ($0{,}3\text{ m}$)}\\[4pt]
  {\scriptsize Enterramento e apiloamento\\da borda superior contra\\escoamento sub-manta}
};

\node[bsTint=bsGold, text width=3.05cm, minimum height=2.75cm, right=of p3] (p4) {%
  \tikz\node[pill=bsGoldDark]{ETAPA 4};\\[6pt]
  \textbf{Desenrolamento\\Montante $\rightarrow$ Jusante}\\[4pt]
  {\scriptsize Aplicação no sentido do\\fluxo com sobreposição\\mínima de $10\text{ cm}$}
};

\node[bsTint=bsForest, text width=3.05cm, minimum height=2.75cm, right=of p4] (p5) {%
  \tikz\node[pill=bsForest]{ETAPA 5};\\[6pt]
  \textbf{Grampeamento em\\Quincôncio \& Inspeção}\\[4pt]
  {\scriptsize Fixação escalonada em U\\e vistoria pós-chuva para\\reparo de sulcos/bordas}
};

\draw[bsArrow, bsNavy] (p1) -- (p2);
\draw[bsArrow, bsTeal] (p2) -- (p3);
\draw[bsArrow, bsEarth] (p3) -- (p4);
\draw[bsArrow, bsGoldDark] (p4) -- (p5);

% Título e subtítulo ancorados acima da faixa
\node[bsSub, above=0.45cm of p3] (sub)
  {Sequência executiva para contato íntimo manta–solo, ancoragem hidráulica e verificação em campo};
\node[bsTitle, above=0.14cm of sub] (title)
  {Procedimento Construtivo e Controle Pós-Instalação de Biomantas};

% Barra de síntese inferior bem separada
\node[fill=bgNavy, draw=bsNavy!40, rounded corners=4pt, inner xsep=10pt, inner ysep=5pt,
      font=\scriptsize\sffamily, text=bsInk, below=0.55cm of p3]
  {\textbf{Ponto crítico de controle:} contato contínuo da manta com o solo — qualquer vão suspenso favorece erosão em túnel sob a biomanta.};

\end{tikzpicture}
\end{document}
